\documentclass[11pt]{article}
\usepackage[T1]{fontenc}
\usepackage{lmodern}
\usepackage[margin=1in]{geometry}
\usepackage{amsmath,amssymb,amsthm}
\usepackage{aliascnt}
\usepackage{microtype}
\usepackage{needspace}
\usepackage[colorlinks=true,linkcolor=blue,citecolor=blue,urlcolor=blue]{hyperref}
\usepackage[nameinlink,noabbrev]{cleveref}
\newtheorem{theorem}{Theorem}[section]
\newaliascnt{proposition}{theorem}
\newtheorem{proposition}[proposition]{Proposition}
\aliascntresetthe{proposition}
\newaliascnt{lemma}{theorem}
\newtheorem{lemma}[lemma]{Lemma}
\aliascntresetthe{lemma}
\theoremstyle{remark}
\newtheorem{remark}[theorem]{Remark}
\newcommand{\R}{\mathbb R}
\newcommand{\dd}{\,\mathrm d}
\newcommand{\E}{\mathcal E}
\newcommand{\vol}{\operatorname{Vol}}
\title{Large Steklov eigenvalues with a prescribed embedded boundary}
\author{Kaiwen Jin\\[2pt]\small\texttt{yui.ui.siki@gmail.com}}
\date{}
\hypersetup{pdftitle={Large Steklov eigenvalues with a prescribed embedded boundary},pdfauthor={Kaiwen Jin},pdfsubject={Geometric spectral theory}}
\begin{document}
\maketitle

\begin{abstract}
For every smooth compact connected embedded boundary of positive dimension,
the first positive Steklov eigenvalue of its smooth embedded fillings is
unbounded, provided the ambient codimension is at least two and a filling
exists. The boundary remains fixed pointwise, and the fillings carry their
induced Euclidean metrics; volume and topology are unrestricted. We give a
detailed proof using intrinsic large-eigenvalue constructions and relative
smooth approximation of short embeddings. This follows the relative
Nash--Kuiper mechanism announced by Girouard and Polymerakis, with the
boundary and arbitrary-curve applications written out explicitly.
\end{abstract}

\section{Introduction}\label{sec:introduction}

The Steklov spectrum measures harmonic functions through their boundary
values. On a smooth compact connected Riemannian manifold $(M,g)$ with
nonempty boundary, it is the spectrum of
\[
 \Delta_g u=0\quad\hbox{in }M,\qquad
 \partial_\nu u=\sigma u\quad\hbox{on }\partial M,
\]
where $\nu$ is the outward unit normal. We write
$0=\sigma_0<\sigma_1\leq\sigma_2\leq\cdots$.
Fixing the boundary metric still leaves considerable freedom in the
interior: in dimensions at least three, conformal changes can make
$\sigma_1$ arbitrarily large while leaving that boundary metric unchanged
\cite{CEG}. Requiring the whole manifold to be embedded in one Euclidean
space, with a prescribed embedded boundary, imposes an additional
geometric constraint.

Colbois--Girouard--Gittins \cite[Theorem~1.1]{CGG} bound the Steklov
eigenvalues of such submanifolds in terms of their fixed boundary and
their volume. They ask in \cite[Remark~1.7]{CGG} whether the first
positive eigenvalue can nevertheless tend to infinity when the boundary
is fixed and the volume is free. Girouard, reporting joint work with
Polymerakis, subsequently announced a relative Nash--Kuiper construction
for transferring intrinsic examples to embedded ones, and described the
case of a circle in $\R^3$ \cite[p.~2075]{GP}. We follow that mechanism.
The purpose of this note is to give a detailed relative approximation
argument and to carry out its application to every admissible smooth
embedded boundary, including arbitrary embedded curves.

\begin{theorem}\label{thm:main}
Let $d\geq1$ and $m\geq d+2$ be integers. Let
$\Sigma^d\subset\R^m$ be a smooth compact connected embedded submanifold
without boundary which bounds at least one smooth compact connected
embedded $(d+1)$-dimensional submanifold. For every $A>0$ there is a
smooth compact connected embedded filling $\Omega\subset\R^m$ such that
\[
 \partial\Omega=\Sigma,\qquad \sigma_1(\Omega)>A,
\]
where $\Omega$ has its induced Euclidean metric.
\end{theorem}

Thus, if $\mathcal F(\Sigma)$ denotes the entire class of these fillings,
$\sup_{\Omega\in\mathcal F(\Sigma)}\sigma_1(\Omega)=\infty$.
The theorem fixes the boundary as an ambient subset, rather than merely
its intrinsic isometry class. It allows both volume and topology to vary.
The volume dependence in the upper bound of \cite{CGG} is consequently
necessary in this setting.

An intrinsic large-eigenvalue metric is generally not the metric induced
by the chosen embedding. We first match it to the induced metric on a
shrinking boundary collar. We then enlarge it conformally in the interior,
making the original embedding short: its induced metric is no larger
than the target metric. This does not lower the Steklov eigenvalue.
Finitely many local primitive-metric additions approximate
the enlarged metric by the metric of a smooth embedding fixed on that
collar. Uniform metric approximation controls the energy quotient and
therefore retains the large eigenvalue. For surfaces, conformal
enlargement leaves this energy unchanged; the intrinsic eigenvalues
instead become large by allowing genus to increase, using the
connected-boundary examples of \cite{CGR}.

Section~\ref{sec:relative} proves the relative approximation and spectral
comparison. Section~\ref{sec:proof} applies them in the two filling
dimensions $n=d+1\geq3$ and $n=2$. All embeddings below take values in
the original ambient $\R^m$, and $\sigma_1$ always means the first
positive eigenvalue.

\section{Relative approximation and the Steklov quotient}\label{sec:relative}

Throughout this section, $M$ is a smooth compact connected manifold of
dimension $n\geq2$ with nonempty boundary. Metrics and maps are smooth up to
the boundary. We use Riemannian volume densities, so orientability is not
assumed. For a metric $g$, put
\[
 \E_g(u)=\int_M |\mathrm du|_g^2\dd V_g.
\]
The first positive Steklov eigenvalue has the characterization
\begin{equation}\label{eq:quotient}
 \sigma_1(M,g)=
 \inf_{\substack{u\in C^\infty(M),\ \int_{\partial M}u\dd S_g=0\\
                         \int_{\partial M}u^2\dd S_g>0}}
 \frac{\E_g(u)}{\int_{\partial M}u^2\dd S_g}.
\end{equation}
This follows from the smooth min--max formula in
\cite[Section~2, equation~(2)]{CEG}. Indeed, on the span of $1$ and a
mean-zero function the maximum quotient is the quotient of that function.
Conversely, every two-dimensional subspace contains a nonzero function
with zero boundary mean. A zero boundary norm is assigned infinite
quotient. The sources \cite{CEG,CGR} number the zero eigenvalue as
$\sigma_1$; their $\sigma_2$ is our $\sigma_1$.

\subsection{Metric comparison}

\begin{lemma}\label{lem:comparison}
Suppose $q$ and $h$ have the same tangential boundary metric and
\[
 (1-t)h\leq q\leq(1+t)h,\qquad 0<t<1.
\]
Then
\begin{equation}\label{eq:comparison}
 \sigma_1(M,q)\geq c_n(t)\sigma_1(M,h),
 \qquad c_n(t)=\frac{(1-t)^{n/2}}{1+t}.
\end{equation}
If $h=\lambda^2k$, where $\lambda\geq1$ and
$\lambda=1$ on $\partial M$, then
\begin{equation}\label{eq:conformal}
 \sigma_1(M,h)\geq\sigma_1(M,k).
\end{equation}
For $n=2$, equality holds in \eqref{eq:conformal}.
\end{lemma}

\begin{proof}
The inverse metrics and volume densities satisfy
\[
 q^{-1}\geq(1+t)^{-1}h^{-1},\qquad
 \dd V_q\geq(1-t)^{n/2}\dd V_h.
\]
Thus $\E_q(u)\geq c_n(t)\E_h(u)$. The boundary measures,
mean-zero condition and denominators in \eqref{eq:quotient} are identical,
which proves \eqref{eq:comparison}. For the second assertion,
\[
 \E_{\lambda^2 k}(u)
   =\int_M\lambda^{n-2}|\mathrm du|_k^2\dd V_k.
\]
The boundary measure is again unchanged. The factor in the integrand is
at least one for $n\geq2$, and is exactly one for $n=2$.
\end{proof}

\subsection{Smooth embeddings fixed near the boundary}

The approximation required below concerns the induced metric uniformly,
and therefore only first derivatives of the embedding. We recall the local
construction that also keeps a boundary collar fixed. A \emph{primitive
metric} on a coordinate chart is a tensor $a^2(\mathrm d\ell)^2$, where
$a$ is smooth with compact support in that chart and $\ell$ is linear in
its coordinates.

For a symmetric tensor $T$ and a positive metric $h$, write
\[
 \|T\|_h=
 \sup_{x\in M}\sup_{0\ne v\in T_xM}
             \frac{|T_x(v,v)|}{h_x(v,v)}.
\]

\begin{lemma}\label{lem:relative-approximation}
Let $f:M^n\to\R^m$ be a smooth embedding, $m>n$, and let
$e=f^*g_{\mathrm E}$ be its induced Euclidean metric. Suppose
\begin{equation}\label{eq:factored-defect}
 h-e=\psi^2B,
\end{equation}
where $B$ is a smooth positive definite metric and $\psi$ is smooth with
compact support in the interior of $M$. For every $\epsilon\in(0,1)$ there
is a smooth embedding $F:M\to\R^m$, equal to $f$ near the whole boundary,
such that
\begin{equation}\label{eq:near-isometry}
 (1-\epsilon)h\leq F^*g_{\mathrm E}\leq(1+\epsilon)h.
\end{equation}
\end{lemma}

\begin{proof}
We use the primitive addition lemma of Eliashberg--Mishachev
\cite[Lemma~21.6.1, pp.~196--197]{EM}. In ambient dimension $m>n$,
an embedding whose metric defect is one primitive metric can be
approximated arbitrarily closely in $C^0$ by embeddings whose induced
metrics approximate the target metric arbitrarily closely. We need the
relative form of this finite construction; it follows from the
parametric one-dimensional lemma used in that proof, as described next.

Work first in an interior coordinate box with primitive defect
$a^2(\mathrm d\ell)^2$, and put
$g=f^*g_{\mathrm E}+a^2(\mathrm d\ell)^2$. The source construction uses the
leaves of $\ker\mathrm d\ell$ as transverse parameters and their
$g$-normal curves as the longitudinal variable. A local Euclidean unit
normal field $\nu$ to the image suffices. If $v$ is the longitudinal
$g$-unit vector, the formal longitudinal derivative
\[
 f_*v+a\,\mathrm d\ell(v)\nu
\]
has Euclidean squared norm
$e(v,v)+a^2\mathrm d\ell(v)^2=g(v,v)=1$. It depends smoothly on the data
and is the genuine derivative wherever $a=0$.
Lemma~17.5.1(c),(d) and the first remark after it in \cite[pp.~161--162]{EM}
keep the construction fixed near the transverse faces and the longitudinal
endpoints when the data there are already genuine. Since $a$ vanishes
near the box boundary, the resulting perturbation is fixed near all its
faces. Refining the support into smaller boxes, if needed, gives the
same statement on an interior chart. This is an application of the
relative parameter clauses to the proof of primitive addition, rather
than an assumption of a smooth exactly isometric embedding.

The local output is an embedding by the same primitive addition lemma.
To retain global injectivity, let $K$ contain the compact support of one
perturbation and choose a slightly larger chart $U$ with $K\subset U$.
For the current embedding, $f(K)$ and $f(M\setminus U)$ are disjoint
compact sets, with positive Euclidean separation. Make the available
$C^0$ error less than half that separation. Points in $K$ cannot meet
the complement of $U$; pairs in $U$ are protected by the local embedding,
and all other pairs are unchanged. The perturbation therefore extends
by $f$ to a global embedding. Coordinate boxes require only local normal
fields, even when the global normal bundle is nontrivial.

We now decompose \eqref{eq:factored-defect}. The positive-cone argument
in \cite[Lemma~21.3.1, p.~193]{EM}, applied to $B$ on compact interior
chart supports, and a smooth square partition of unity give
\begin{equation}\label{eq:primitive-sum}
 \psi^2 B=\sum_{j=1}^N p_j,
 \qquad p_j=a_j^2(\mathrm d\ell_j)^2.
\end{equation}
Here every $a_j$ is smooth and supported compactly in one interior chart.
For clarity, the square amplitudes can be obtained by first writing the
positive matrix $B(x)$, locally, as a positive smooth combination of
finitely many constant positive matrices. Choose these matrices as
vertices of a small polytope containing the compact range under
consideration, and split each vertex into squares of linear forms.
The coefficients are strictly positive on smaller chart supports.
Multiplication by $\psi$ and the square-partition functions then gives
$a_j=\psi\rho_j\sqrt{\beta_j}$ with smooth positive $\beta_j$.
There is no square root of a vanishing arbitrary coefficient to regularize.

Starting from $f_0=f$, at step $j$ use the target
\[
 g_j=f_{j-1}^*g_{\mathrm E}+p_j.
\]
The initial map at that step is short for $g_j$ exactly. Choose the
new metric error
$E_j=f_j^*g_{\mathrm E}-g_j$ so that $\|E_j\|_h<\epsilon/N$,
along with the $C^0$ tolerance just described. This is possible because
$g_j\leq C_jh$ for some finite $C_j$ on compact $M$, so relative
$g_j$ error less than $\epsilon/(NC_j)$ suffices. The finite construction
telescopes to
\[
 f_N^*g_{\mathrm E}
   =e+\sum_{j=1}^Np_j+\sum_{j=1}^NE_j
   =h+\sum_{j=1}^NE_j,
\]
which proves \eqref{eq:near-isometry}.

All modifications are supported away from $\partial M$, and the
intersection of their finitely many unchanged collars is still a
boundary collar. The finite smooth constructions above give a smooth
map. One may also use a $C^1$ formulation of the primitive lemma and
then smooth relative to that collar with arbitrarily small $C^1$ error.
Embeddings of a compact manifold are open in $C^1$: local graph estimates
preserve injectivity on finitely many small charts, and compact separation
handles the remaining pairs. Such smoothing preserves the embedding and
changes its induced metric arbitrarily little. Reserving part of each
error tolerance for smoothing gives the same conclusion.
\end{proof}

\subsection{Passing from an intrinsic metric to an embedded filling}

The intrinsic metric need not initially dominate the induced metric.
A conformal enlargement can arrange this without decreasing the
Steklov quotient. We first make it agree with the induced metric on a
small collar.

\Needspace{8\baselineskip}
\begin{proposition}\label{prop:transfer}
Let $f:M^n\to\R^m$ be a smooth embedding, $n\geq2$, $m>n$, and
let $k$ be any smooth metric whose tangential boundary restriction
equals that of $e=f^*g_{\mathrm E}$. For every $\theta\in(0,1)$ there is
a smooth embedding $F:M\to\R^m$, equal to $f$ near $\partial M$, with
\begin{equation}\label{eq:transfer}
 \sigma_1(F(M))\geq(1-\theta)\sigma_1(M,k).
\end{equation}
The image carries its induced Euclidean metric.
\end{proposition}

\begin{proof}
Equality of tangential metrics does not match their normal components.
Let $C_e$ and $C_k$ be collars obtained by flowing extensions of the
inward unit normal fields of $e$ and $k$. At the boundary,
$C_k\circ C_e^{-1}$ fixes tangent vectors and sends the $e$-unit inward
normal to the $k$-unit inward normal. It therefore pulls $k$ back to $e$
as a full tensor on $TM|_{\partial M}$.

This collar map extends to a boundary-fixing diffeomorphism of $M$.
Indeed, interpolate the two inward transverse fields linearly. Their
flows give an isotopy of collars on a common short interval, by
compactness. On the image of a smaller collar its velocity in the
isotopy parameter is a smooth time-dependent vector field, zero at
boundary points. Cut it off inside a larger collar and integrate.
Its flow agrees with the collar isotopy near the boundary and is
the identity away from the larger collar. These flow and collar facts
are the usual smooth constructions; see \cite[Section~6.2]{Hirsch}.
Pulling back $k$ by this diffeomorphism preserves its Steklov spectrum.
Thus we may assume
\begin{equation}\label{eq:full-boundary}
 k=e\quad\hbox{on }TM|_{\partial M}.
\end{equation}

In an $e$-normal flow collar, let $t\geq0$ be the inward coordinate.
Choose $0\leq\psi_\delta\leq1$, smooth, equal to zero for $t\leq\delta$
and one for $t\geq2\delta$, extended as one over the remainder of $M$.
Set $\chi=\psi_\delta^2$ and
\[
 k_\delta=(1-\chi)e+\chi k.
\]
By \eqref{eq:full-boundary}, the uniform $e$-norm of $k-e$ is $O(t)$
in the collar. Since $k_\delta-k$ is supported where $t<2\delta$,
$k_\delta\to k$ uniformly. For any fixed $\eta\in(0,1)$ we can choose
$\delta>0$ so that
\[
 (1-\eta)k\leq k_\delta\leq(1+\eta)k,
 \qquad
 \sigma_1(M,k_\delta)\geq c_n(\eta)\sigma_1(M,k).
\]
The last inequality is \cref{lem:comparison}. This collar width is
chosen separately for the given $k$.

Choose $c\in(0,1]$ such that $k\geq ce$ on $M$, and then
$K>(1-c)/c$. Put
\[
 h=(1+K\chi)k_\delta.
\]
The conformal factor is at least one and equals one on the boundary
collar. Thus $\sigma_1(M,h)\geq\sigma_1(M,k_\delta)$, including
dimension two. Moreover,
\begin{align*}
 h-e&=\chi B,\\
 B&=Ke+(1+K\chi)(k-e)
    \geq\bigl[Kc-(1-c)\bigr]e>0.
\end{align*}
The displayed formula defines a smooth positive $B$ even at $\chi=0$.
Consequently $h-e=\psi_\delta^2 B$ has the form required by
\cref{lem:relative-approximation}.

Apply that lemma with tolerance $\epsilon$. The induced metric
$q=F^*g_{\mathrm E}$ satisfies \eqref{eq:near-isometry} and has exactly
the boundary metric of $h$. Hence
\[
 \sigma_1(F(M))
 \geq c_n(\epsilon)\sigma_1(M,h)
 \geq c_n(\epsilon)c_n(\eta)\sigma_1(M,k).
\]
For this fixed $n$, choose $\eta,\epsilon$ small enough that the product
exceeds $1-\theta$. This proves \eqref{eq:transfer}.
\end{proof}

\section{Proof of the prescribed-boundary theorem}\label{sec:proof}

Put $n=d+1$. We use \cref{prop:transfer} first on a fixed filling in
dimensions $n\geq3$, and then on a family of embedded surfaces when
$n=2$.

\subsection{Filling dimension at least three}

Choose any smooth compact connected embedded filling $M$ of $\Sigma$
and let $e$ be its induced Euclidean metric. Colbois--El Soufi--Girouard
\cite[Theorem~1.1(ii)]{CEG} prove that on every compact connected
Riemannian manifold of dimension at least three with nonempty boundary,
there are smooth conformal metrics equal to the original metric on the
boundary whose first positive Steklov eigenvalues tend to infinity.
In particular, for $A>0$ choose such a metric $k$ with
$\sigma_1(M,k)>2A$.
Applying \cref{prop:transfer} to the inclusion $M\hookrightarrow\R^m$
with $\theta=1/2$ gives a smooth compact connected embedded image
$\Omega$ with $\partial\Omega=\Sigma$ and $\sigma_1(\Omega)>A$.
No orientability assumption on $M$ enters the argument.

\subsection{Surface fillings of an arbitrary embedded curve}

Now $d=1$ and $m\geq3$. Since $\Sigma$ is a compact connected smooth
one-manifold without boundary, it is a circle. We first construct
orientable embedded fillings of every sufficiently large genus.

When $m=3$, Seifert's construction gives an orientable spanning surface
for the knot $\Sigma$ \cite[Section~1, pp.~572--573]{Seifert}.
To pass from its polygonal formulation to the smooth original curve,
take a sufficiently fine polygon in a normal tubular neighborhood of
$\Sigma$ and smooth its vertices. The resulting curve is a normal graph
over $\Sigma$, and interpolation of its normal displacement is a smooth
isotopy. The finite disks and bands in Seifert's construction can be
smoothed. Extending the knot isotopy to the ambient space transports
this surface to one whose boundary is exactly the original $\Sigma$;
see \cite[Theorems~4.5.1 and~8.1.3]{Hirsch} for tubular neighborhoods
and isotopy extension.

For $m>3$ there is a rank-three linear map
$\pi:\R^m\to\R^3$ which embeds $\Sigma$. Here is the general-position
argument. On the open space of rank-three linear maps, the evaluation
\[
 (P,x,y)\longmapsto P(x-y),\qquad x,y\in\Sigma,\quad x\ne y,
\]
is a submersion, since $x-y\ne0$. Parametric transversality
\cite[Theorem~3.2.7]{Hirsch} gives a residual set of $P$ for which its
restriction to the two-dimensional pair space is transverse to zero
in $\R^3$, and therefore has no zero. The same argument on the
one-dimensional unit tangent bundle of $\Sigma$, using $(P,x,v)\mapsto
P(v)$, gives immersion. The intersection of these residual sets is
nonempty; a compact injective immersion is an embedding.

Let $\Gamma=\pi(\Sigma)$, and choose a smooth orientable Seifert surface
$S$ for $\Gamma$ in $\R^3$. A linear splitting of $\pi$ identifies the
ambient vector space with $\R^3\times\R^{m-3}$, in which $\Sigma$ is the
graph of a smooth function $z:\Gamma\to\R^{m-3}$. Extend $z$ over $S$
by multiplying its collar extension by a cutoff which vanishes in the
interior. The graph over $S$ of this extension is embedded and has
boundary $\Sigma$. Inverting the linear splitting gives a smooth
orientable filling in the original $\R^m$. The splitting is used only
for this topological construction; all induced metrics are taken in
the original Euclidean coordinates.

Let $S_0$ be an orientable filling obtained in this way, of genus $g_0$.
In a ball around an interior disk, disjoint from the rest of the surface
and its boundary, replace a smaller disk by a punctured torus joined
along a thin annulus. In local ambient coordinates this is the usual
embedded connected sum in a three-dimensional slice, available for
every $m\geq3$. It is smooth after rounding the seams, preserves
orientability and the outer disk collar, and decreases Euler
characteristic by two. Repetition gives an embedded orientable filling
$M_g$ of genus $g$ for every $g\geq g_0$, with boundary exactly $\Sigma$.

Colbois--Girouard--Raveendran \cite[Theorem~7]{CGR} construct compact
surfaces $(S_l,k_l)$ of genus $1+l$, with connected boundary, such that
\begin{equation}\label{eq:surface-source}
 \sigma_1(S_l,k_l)\,\operatorname{length}_{k_l}(\partial S_l)
       \geq Cl,
\end{equation}
where $C>0$ is independent of $l$. Their construction can be chosen
orientable. Section~6.4 starts from a connected four-regular expander
graph $G_l$, placing a genus-zero block with four identical product
boundary collars at each vertex and pairing circles along graph edges.
Orientation-reversing identifications of paired circles give an
orientable closed surface. A maximal tree in $G_l$ is embedded in this
surface by joining points in the blocks through the corresponding
collars. Its regular neighborhood is a disk; removing it leaves one
boundary component. To keep the geometry uniform, the source implements
this removal with the four carved piece types in its Figure~5, selected
according to the vertex degrees $1,2,3,4$ in the maximal tree.
Orientations can be chosen along
the tree first, and the remaining full circles glued by
orientation-reversing collar isometries. These choices leave the pieces
and their metric estimates unchanged; the freedom in gluing and its
irrelevance to the spectral conclusion are discussed in
\cite[Remark~37]{CGR}.

Let $L=\operatorname{length}_{g_{\mathrm E}}(\Sigma)>0$ and
$L_l=\operatorname{length}_{k_l}(\partial S_l)$. Rescale the intrinsic
metrics by
\[
 k'_l=(L/L_l)^2k_l.
\]
For a constant $a>0$, boundary length scales by $a$ and Steklov
eigenvalues by $a^{-1}$ under $k\mapsto a^2k$. Thus
\begin{equation}\label{eq:surface-scaled}
 \operatorname{length}_{k'_l}(\partial S_l)=L,
 \qquad \sigma_1(S_l,k'_l)\geq Cl/L.
\end{equation}

For sufficiently large $l$, the embedded filling $M_{1+l}$ exists and
is diffeomorphic to $S_l$, since both are compact connected orientable
surfaces with the same genus and one boundary component
\cite[Theorem~9.3.7]{Hirsch}. The boundary value of the diffeomorphism
may be chosen to be an arclength isometry. To see this explicitly,
start with a surface diffeomorphism $T:M_{1+l}\to S_l$, choose compatible
orientations, and let $b:\Sigma\to\partial S_l$ be an orientation-compatible
arclength isometry for $k'_l$. The discrepancy
$(T|_\Sigma)^{-1}\circ b$ is an orientation-preserving circle
diffeomorphism. In the arclength coordinate of period $L$, lift it to
a smooth increasing function $\beta$ with
$\beta(s+L)=\beta(s)+L$. The functions
\[
 \beta_t(s)=(1-t)s+t\beta(s),\qquad 0\leq t\leq1,
\]
have positive derivative. On a boundary collar use
$(s,r)\mapsto(\beta_{\tau(r)}(s),r)$, where $\tau=1$ near the boundary
and $\tau=0$ near the inner edge. This is a diffeomorphism, extends by
identity, and adjusts $T$ to have boundary value $b$.

Pull $k'_l$ back by this adjusted diffeomorphism to a metric $k$ on
the actual embedded filling $M_{1+l}$. Its tangential boundary metric
is exactly the Euclidean metric on $\Sigma$. By
\eqref{eq:surface-scaled}, choose $l$ so that $\sigma_1(M_{1+l},k)>2A$.
Apply \cref{prop:transfer} to the inclusion of this surface with
$\theta=1/2$. Its smooth embedded image $\Omega$ has
$\partial\Omega=\Sigma$ and $\sigma_1(\Omega)>A$.
This proves the theorem in the remaining dimension and hence for all
the stated $d,m,\Sigma$ and $A$.\qed

\section*{AI-use statement}
OpenAI's Codex assisted with the mathematical argument, source checking,
and manuscript preparation. An earlier draft was examined in separate
fresh Codex contexts for reader comprehension and referee-style
objections. Codex then assisted with revision and final checks.
No independent human review or formal verification is claimed.

\begin{thebibliography}{9}

\bibitem{CEG}
B.~Colbois, A.~El Soufi and A.~Girouard,
\emph{Compact manifolds with fixed boundary and large Steklov eigenvalues},
\href{https://arxiv.org/abs/1701.04125v2}{arXiv:1701.04125v2} (2018).

\bibitem{CGG}
B.~Colbois, A.~Girouard and K.~Gittins,
\emph{Steklov eigenvalues of submanifolds with prescribed boundary in
Euclidean space},
\href{https://arxiv.org/abs/1711.06458v2}{arXiv:1711.06458v2} (2018).

\bibitem{CGR}
B.~Colbois, A.~Girouard and B.~Raveendran,
\emph{The Steklov spectrum and coarse discretizations of manifolds with
boundary},
\href{https://arxiv.org/abs/1612.07665v1}{arXiv:1612.07665v1} (2016).

\bibitem{EM}
Y.~Eliashberg and N.~Mishachev,
\emph{Introduction to the $h$-Principle},
Graduate Studies in Mathematics~48, American Mathematical Society, 2002.
\href{https://doi.org/10.1090/gsm/048}{doi:10.1090/gsm/048}.
Page references to this work use the
\href{https://virtualmath1.stanford.edu/~eliash/Public/257B-2020/intro.pdf}{author-hosted text}.

\bibitem{GP}
A.~Girouard,
\emph{Manifolds with arbitrarily large Steklov eigenvalues}
(joint work with P.~Polymerakis), in \emph{Geometric Spectral Theory},
Oberwolfach Report~36/2023, pp.~2074--2076.
\href{https://doi.org/10.4171/OWR/2023/36}{doi:10.4171/OWR/2023/36}.

\bibitem{Hirsch}
M.~W.~Hirsch,
\emph{Differential Topology},
Graduate Texts in Mathematics~33, Springer-Verlag, 1976.
\href{https://doi.org/10.1007/978-1-4684-9449-5}{doi:10.1007/978-1-4684-9449-5}.

\bibitem{Seifert}
H.~Seifert,
\emph{\"Uber das Geschlecht von Knoten},
Mathematische Annalen~110 (1935), 571--592.
\href{https://doi.org/10.1007/BF01448044}{doi:10.1007/BF01448044}.

\end{thebibliography}
\end{document}
